\renewcommand{\deutschertitel}{TBD}
\renewcommand{\englischertitel}{Companion Computer}
\chapter[\texorpdfstring{\protect{\vspace{2pt}\englischertitel}}{\englischertitel}]{}
\kapitel{\deutschertitel}
\thispagestyle{empty}

\label{KapitelCompanionComputer}

\begin{paracol}{2}[]

{\raggedright\huge\bfseries\sffamily \englischertitel \par\ } \\[1.8ex]

\switchcolumn

{\raggedright\huge\bfseries\sffamily \deutschertitel \par\ } \\[1.8ex]

\coleng

A companion computer is useful for handling computationally intensive tasks such as artificial intelligence, computer vision, dynamic path replanning, cloud communication, and mission planning, while leaving the drone’s stabilization and low-level flight control to the flight controller. Common examples of companion computers include the Raspberry Pi and NVIDIA Jetson platforms.

\colger

TBD

\colende

\renewcommand{\deutschertitel}{TBD}
\renewcommand{\englischertitel}{Raspberry Pi as a Companion Computer}
\makroabschnitt
\label{AbschnittRaspberryPiCompanionComputer}

There are different versions of the Raspberry Pi, each serving different purposes depending on performance requirements and intended applications. To assist in selecting the appropriate model, a brief description of the available versions is provided in Section~\ref{AbschnittObjekterkennungHuckepack}.

\colger

TBD

\colende

\renewcommand{\deutschertitel}{TBD}
\renewcommand{\englischertitel}{Flashing Raspberry Pi OS}
\makrounterabschnitt

It is recommended to install the headless version of Raspberry Pi OS, which is designed to operate without a monitor, keyboard, or mouse. Since the Raspberry Pi serves as a companion computer mounted on the UAV and is intended to be used remotely, a headless operating system is the most suitable choice. In addition, the absence of a graphical desktop environment reduces CPU utilization, memory consumption, and power usage.

\colger

TBD

\coleng

To flash Raspberry Pi OS, the following components are required:

\colger

TBD

\coleng

\begin{itemize}
	\item Raspberry pi
\end{itemize}

\begin{itemize}
	\item MicroSD card (at least 4 GB)
\end{itemize}

\begin{itemize}
	\item MicroSD card reader
\end{itemize}

\begin{itemize}
	\item Raspberry Pi Imager software (available for Windows, macOS, and Linux)
\end{itemize}

\colger

TBD

\colende

\begin{figure}[htbp]
  \centering
  \begin{minipage}[t]{0.48\textwidth}
    \centering
    \includegraphics[width=\linewidth]{pi_imager_screenshot_1.png}
    \vspace{0pt} % sorgt für Top-Ausrichtung
  \end{minipage}\hfill
  \begin{minipage}[t]{0.48\textwidth}
    \centering
    \includegraphics[width=\linewidth]{pi_imager_screenshot_2.png}
    \vspace{0pt}
  \end{minipage}
  \caption{Raspberry Pi Imager Software Screenshots. Left: OS selection. Right: Network configuration.}
  \label{AbbildungRaspberryPiImagerSoftwareScreenshots1}
\end{figure}
\FloatBarrier

\colstart

Next, open Raspberry Pi Imager and select the appropriate Raspberry Pi model. Under Operating System, choose Raspberry Pi OS (Other) → Raspberry Pi OS Lite (See Figure~\ref{AbbildungRaspberryPiImagerSoftwareScreenshots1} left).

\colger

TBD

\coleng

Before writing the operating system to the microSD card, it is essential to do some configuration.

\colger

TBD

\coleng

\begin{itemize}
	\item Configure the home WiFi network by entering the SSID and WiFi password (See Figure~\ref{AbbildungRaspberryPiImagerSoftwareScreenshots1} right).
\end{itemize}

\begin{itemize}
	\item Configure the username and password to access the Raspberry Pi remotely, or enable SSH access with a public key (See Figure~\ref{AbbildungRaspberryPiImagerSoftwareScreenshots2} left).
\end{itemize}

\colende

\begin{figure}[htbp]
  \centering
  \begin{minipage}[t]{0.48\textwidth}
    \centering
    \includegraphics[width=\linewidth]{pi_imager_screenshot_3.png}
    \vspace{0pt} % sorgt für Top-Ausrichtung
  \end{minipage}\hfill
  \begin{minipage}[t]{0.48\textwidth}
    \centering
    \includegraphics[width=\linewidth]{pi_imager_screenshot_4.png}
    \vspace{0pt}
  \end{minipage}
  \caption{Raspberry Pi Imager Software Screenshots. Left: Password configuration. Right: Configuration summary.}
  \label{AbbildungRaspberryPiImagerSoftwareScreenshots2}
\end{figure}
\FloatBarrier

\colstart

Once the configuration is complete, click Write to flash the operating system onto the microSD card. After the flashing process has finished, insert the microSD card into the Raspberry Pi and power on the device.

\colger

TBD

\colende

\renewcommand{\deutschertitel}{TBD}
\renewcommand{\englischertitel}{Configuring Remote Communication}
\makrounterabschnitt

After booting, the Raspberry Pi automatically attempts to connect to the configured home WiFi network. If the connection is successful, the router assigns an IP address to the device.

\colger

TBD

\coleng

To determine the assigned IP address, open a command prompt on a computer connected to the same network and execute:

\colger

TBD

\colende

\lstdefinestyle{shell}{
  backgroundcolor=\color{gray!10},
  basicstyle=\ttfamily\small,
  frame=single,
  breaklines=true,        % <-- Zeilenumbruch aktivieren
  breakatwhitespace=true, % <-- nur bei Leerzeichen umbrechen 
  postbreak=\mbox{\textcolor{gray}{$\hookrightarrow$}\space}, % Pfeil am Zeilenende
  showstringspaces=false,
  xleftmargin=0em,
  xrightmargin=0em,
  framerule=0.5pt,
  rulecolor=\color{gray!60}
}

\begin{lstlisting}[style=shell]
ping <hostname>.local
//example: ping raspberrypi.local
\end{lstlisting}

\colstart

Alternatively, the IP address can be obtained from the list of connected devices in the router's administration interface.

\colger

TBD

\coleng

Once the IP address has been identified, establish an SSH connection using:

\colger

TBD

\colende

\lstdefinestyle{shell}{
  backgroundcolor=\color{gray!10},
  basicstyle=\ttfamily\small,
  frame=single,
  breaklines=true,        % <-- Zeilenumbruch aktivieren
  breakatwhitespace=true, % <-- nur bei Leerzeichen umbrechen 
  postbreak=\mbox{\textcolor{gray}{$\hookrightarrow$}\space}, % Pfeil am Zeilenende
  showstringspaces=false,
  xleftmargin=0em,
  xrightmargin=0em,
  framerule=0.5pt,
  rulecolor=\color{gray!60}
}

\begin{lstlisting}[style=shell]
ssh <username>@<ip_address>
//example: ssh pi@192.168.0.153
\end{lstlisting}

\renewcommand{\deutschertitel}{TBD}
\renewcommand{\englischertitel}{Configuring a Fallback WiFi Hotspot}
\makrounterunterabschnitt

For field operations, a known WiFi network may not always be available. Therefore, it is useful to configure the Raspberry Pi so that it automatically creates its own wireless access point whenever it cannot connect to a predefined network. The desired behavior is as follows:

\colger

TBD

\coleng

\begin{itemize}
	\item Known WiFi network available: Connect to the known network, allwing the other devices on the network to communicate with the Raspberry Pi.
\end{itemize}

\begin{itemize}
	\item \item No known WiFi network available: Create a wireless access point for the Raspberry Pi, enabling direct communication with other devices without relying on an external network.
\end{itemize}

\colger

TBD

\colende

\renewcommand{\deutschertitel}{TBD}
\renewcommand{\englischertitel}{Step 1: Create a new WiFi connection using NetworkManager:}
\makrounterunterabschnitt

\colende

\lstdefinestyle{shell}{
  backgroundcolor=\color{gray!10},
  basicstyle=\ttfamily\small,
  frame=single,
  breaklines=true,        % <-- Zeilenumbruch aktivieren
  breakatwhitespace=true, % <-- nur bei Leerzeichen umbrechen 
  postbreak=\mbox{\textcolor{gray}{$\hookrightarrow$}\space}, % Pfeil am Zeilenende
  showstringspaces=false,
  xleftmargin=0em,
  xrightmargin=0em,
  framerule=0.5pt,
  rulecolor=\color{gray!60}
}

\begin{lstlisting}[style=shell]
sudo nmcli connection add \
  type wifi \
  ifname wlan0 \
  con-name Pi-Hotspot \
  autoconnect no \
  ssid RespberryPi-Hotspot \
\end{lstlisting}

\colstart

The connection should be configured as an access point:

\colger

TBD

\colende

\lstdefinestyle{shell}{
  backgroundcolor=\color{gray!10},
  basicstyle=\ttfamily\small,
  frame=single,
  breaklines=true,        % <-- Zeilenumbruch aktivieren
  breakatwhitespace=true, % <-- nur bei Leerzeichen umbrechen 
  postbreak=\mbox{\textcolor{gray}{$\hookrightarrow$}\space}, % Pfeil am Zeilenende
  showstringspaces=false,
  xleftmargin=0em,
  xrightmargin=0em,
  framerule=0.5pt,
  rulecolor=\color{gray!60}
}

\begin{lstlisting}[style=shell]
sudo nmcli connection modify Pi-Hotspot \
    802-11-wireless.mode ap \
    ipv4.method shared \
    wifi-sec.key-mgmt wpa-psk \
    wifi-sec.psk "ChooseAStrongPassword"
\end{lstlisting}

\colstart

To verify that both the home WiFi connection and the hotspot configuration are stored, execute:

\colger

TBD

\colende

\lstdefinestyle{shell}{
  backgroundcolor=\color{gray!10},
  basicstyle=\ttfamily\small,
  frame=single,
  breaklines=true,        % <-- Zeilenumbruch aktivieren
  breakatwhitespace=true, % <-- nur bei Leerzeichen umbrechen 
  postbreak=\mbox{\textcolor{gray}{$\hookrightarrow$}\space}, % Pfeil am Zeilenende
  showstringspaces=false,
  xleftmargin=0em,
  xrightmargin=0em,
  framerule=0.5pt,
  rulecolor=\color{gray!60}
}

\begin{lstlisting}[style=shell]
nmcli connection show
// This should show:
// HomeWiFi
// Pi-Hotspot
\end{lstlisting}

\renewcommand{\deutschertitel}{TBD}
\renewcommand{\englischertitel}{Step 2: Create a Fallback Script}
\makrounterunterabschnitt

This script checks whether the Raspberry Pi has successfully connected to a WiFi network. If no connection is detected, the hotspot is activated. The script can be saved as \verb!/usr/local/bin/wifi-fallback.sh! using the following command:

\colger

TBD

\colende

\lstdefinestyle{shell}{
  backgroundcolor=\color{gray!10},
  basicstyle=\ttfamily\small,
  frame=single,
  breaklines=true,        % <-- Zeilenumbruch aktivieren
  breakatwhitespace=true, % <-- nur bei Leerzeichen umbrechen 
  postbreak=\mbox{\textcolor{gray}{$\hookrightarrow$}\space}, % Pfeil am Zeilenende
  showstringspaces=false,
  xleftmargin=0em,
  xrightmargin=0em,
  framerule=0.5pt,
  rulecolor=\color{gray!60}
}

\begin{lstlisting}[style=shell]
sudo nano /usr/local/bin/wifi-fallback.sh
\end{lstlisting}

\colstart

Insert the following configuration into the file:

\colger

TBD

\colende

\lstdefinestyle{shell}{
  backgroundcolor=\color{gray!10},
  basicstyle=\ttfamily\small,
  frame=single,
  breaklines=true,        % <-- Zeilenumbruch aktivieren
  breakatwhitespace=true, % <-- nur bei Leerzeichen umbrechen 
  postbreak=\mbox{\textcolor{gray}{$\hookrightarrow$}\space}, % Pfeil am Zeilenende
  showstringspaces=false,
  xleftmargin=0em,
  xrightmargin=0em,
  framerule=0.5pt,
  rulecolor=\color{gray!60}
}

\begin{lstlisting}[style=shell]
#!/bin/bash

# Wait a little after boot
sleep 30

# Check whether WiFi is connected
if nmcli -t -f GENERAL.STATE device show wlan0 | grep -q "100"; then
    exit 0
fi

echo "Could not connect to WiFi. Starting hotspot..."

nmcli connection up Pi-Hotspot
\end{lstlisting}

\colstart

Once the script is saved, make it executable with the following command:

\colger

TBD

\colende

\lstdefinestyle{shell}{
  backgroundcolor=\color{gray!10},
  basicstyle=\ttfamily\small,
  frame=single,
  breaklines=true,        % <-- Zeilenumbruch aktivieren
  breakatwhitespace=true, % <-- nur bei Leerzeichen umbrechen 
  postbreak=\mbox{\textcolor{gray}{$\hookrightarrow$}\space}, % Pfeil am Zeilenende
  showstringspaces=false,
  xleftmargin=0em,
  xrightmargin=0em,
  framerule=0.5pt,
  rulecolor=\color{gray!60}
}

\begin{lstlisting}[style=shell]
sudo nano /etc/systemd/system/wifi-fallback.service
\end{lstlisting}

\colstart

Now, in order to execute the fallback script automatically during system startup, create a systemd service:

\colger

TBD

\colende

\lstdefinestyle{shell}{
  backgroundcolor=\color{gray!10},
  basicstyle=\ttfamily\small,
  frame=single,
  breaklines=true,        % <-- Zeilenumbruch aktivieren
  breakatwhitespace=true, % <-- nur bei Leerzeichen umbrechen 
  postbreak=\mbox{\textcolor{gray}{$\hookrightarrow$}\space}, % Pfeil am Zeilenende
  showstringspaces=false,
  xleftmargin=0em,
  xrightmargin=0em,
  framerule=0.5pt,
  rulecolor=\color{gray!60}
}

\begin{lstlisting}[style=shell]
sudo nano /etc/systemd/system/wifi-fallback.service
\end{lstlisting}

\colstart

Insert the following configuration into the file:

\colger

TBD

\colende

\lstdefinestyle{shell}{
  backgroundcolor=\color{gray!10},
  basicstyle=\ttfamily\small,
  frame=single,
  breaklines=true,        % <-- Zeilenumbruch aktivieren
  breakatwhitespace=true, % <-- nur bei Leerzeichen umbrechen 
  postbreak=\mbox{\textcolor{gray}{$\hookrightarrow$}\space}, % Pfeil am Zeilenende
  showstringspaces=false,
  xleftmargin=0em,
  xrightmargin=0em,
  framerule=0.5pt,
  rulecolor=\color{gray!60}
}

\begin{lstlisting}[style=shell]
[Unit]
Description=Fallback WiFi Hotspot
After=NetworkManager.service

[Service]
Type=oneshot
ExecStart=/usr/local/bin/wifi-fallback.sh

[Install]
WantedBy=multi-user.target
\end{lstlisting}

\colstart

Save the file and enable the service to start automatically during system boot:

\colger

TBD

\colende

\lstdefinestyle{shell}{
  backgroundcolor=\color{gray!10},
  basicstyle=\ttfamily\small,
  frame=single,
  breaklines=true,        % <-- Zeilenumbruch aktivieren
  breakatwhitespace=true, % <-- nur bei Leerzeichen umbrechen 
  postbreak=\mbox{\textcolor{gray}{$\hookrightarrow$}\space}, % Pfeil am Zeilenende
  showstringspaces=false,
  xleftmargin=0em,
  xrightmargin=0em,
  framerule=0.5pt,
  rulecolor=\color{gray!60}
}

\begin{lstlisting}[style=shell]
sudo systemctl daemon-reload
sudo systemctl enable wifi-fallback.service
\end{lstlisting}

\renewcommand{\deutschertitel}{TBD}
\renewcommand{\englischertitel}{Step 3: Verifying the Fallback Configuration}
\makrounterunterabschnitt

To test the fallback mechanism, disable the home WiFi network and reboot the Raspberry Pi. After some time, the Raspberry Pi should create its wireless access point. Connect another device to the hotspot using the configured SSID and password. The Raspberry Pi typically assigns itself the IP address 10.42.0.1, allowing SSH access via:

\colger

TBD

\colende

\lstdefinestyle{shell}{
  backgroundcolor=\color{gray!10},
  basicstyle=\ttfamily\small,
  frame=single,
  breaklines=true,        % <-- Zeilenumbruch aktivieren
  breakatwhitespace=true, % <-- nur bei Leerzeichen umbrechen 
  postbreak=\mbox{\textcolor{gray}{$\hookrightarrow$}\space}, % Pfeil am Zeilenende
  showstringspaces=false,
  xleftmargin=0em,
  xrightmargin=0em,
  framerule=0.5pt,
  rulecolor=\color{gray!60}
}

\begin{lstlisting}[style=shell]
ssh pi@10.42.0.1
// Here pi is the configured username.
\end{lstlisting}

\colstart

TBD

\colger

TBD

\colende

\begin{figure}[htb!]
	\centering
	\includegraphics[width=\linewidth]{Complete_Drone_Wiring_Diagram_with_Raspberry_Pi_Zero_2_W_Companion_Computer.png}
	\caption{Complete Drone Wiring Diagram with Raspberry Pi Zero 2 W Companion Computer.}
	\label{AbbildungCompleteDroneWiringDiagramWithRaspberryPiZero2WCompanionComputer}
\end{figure}
\FloatBarrier
